\chapter{Espacios cociente y teoremas de isomorfía}
\label{cap:cocientes-isomorfia}

El cociente identifica vectores que difieren por un elemento de un subespacio. Esta construcción abstrae la idea de eliminar direcciones redundantes y conduce al teorema de isomorfía para transformaciones lineales \cite{lang1987,hoffmankunze1971,axler2024}.

\section{Clases laterales de un subespacio}
Clases de equivalencia determinadas por un subespacio y su interpretación.

\section{Definición del espacio cociente \texorpdfstring{$V/W$}{V/W}}
Estructura vectorial sobre el conjunto de clases laterales.

\section{Operaciones en espacios cociente}
Suma y multiplicación por escalares bien definidas.

\section{Dimensión del cociente}
Relación entre las dimensiones de $V$, $W$ y $V/W$.

\section{Proyección canónica}
Aplicación lineal de $V$ en $V/W$ y descripción de su núcleo.

\section{Transformaciones inducidas}
Condiciones para que una transformación descienda al cociente.

\section{Primer teorema de isomorfía}
Construcción y demostración del isomorfismo asociado a una transformación lineal.

\section{El cociente por el núcleo}
Isomorfismo natural $V/\ker T\cong\operatorname{im}T$ y su relación con rango--nulidad.
